\scanpage{38}
On procède de $=$ pour $f$, en notant \inserted{comme dessus} qu'il ne peut exister qu'\uline{une seule $f$} : \deleted{\unsure{d'ordre|et de}} \inserted{$P^1_S\mathbin{\unsure{-|\to}}P^1_S$} définissant un \unsure{revêtem\textsuperscript{t}|revêtement} d'ordre donné (ici 3) et dont la valeur en une section (p.ex. 0) disjointe de la ramification -- \unsure{convertie|consistant} en trois sections disjointes -- soit donnée. On trouve, en ramenant \unsure{cette|la} ramification \unsure{à|:} être $0,\infty$ et \unsure{revenant|arrivant|aucune} en arrière

\begingroup\small
\[
\begin{aligned}
s=f(t)&=\frac{1+\left(\frac{\deleted{\unsure{x}}t-x'}{\deleted{\unsure{x'}}t-x''}\right)^3}
{1-\left(\frac{\deleted{\unsure{x}}t-x'}{\deleted{\unsure{x'}}t-x''}\right)^3}
=\frac{(t-x'')^3+(t-x')^3}{(t-x'')^3-(t-x')^3}\\
&=\frac{2t^3-3t^2({x''}^{\deleted{\unsure{2}}}+{x'}^{\deleted{\unsure{2}}})+3t^{\deleted{\unsure{2}}}({x'}^{2}+{x''}^{2})-({x'}^{3}+{x''}^{3})}
{3t^2(x'-x'')+3t({x''}^{2}-{x'}^{2})-({x''}^{3}-{x'}^{3})}
\end{aligned}
\]
\endgroup

\noindent\begin{minipage}[t]{.48\linewidth}\scriptsize
\tikz[overlay,baseline,x=\linewidth,y=1ex]{\draw[thin] (.06,1.8) .. controls (.31,1.6) and (.55,1.5) .. (.69,1.3) .. controls (.73,1.1) and (.74,.2) .. (.76,-.7) .. controls (.8,-1.1) and (.86,-.85) .. (.88,-2.65) .. controls (.88,-3.3) and (.93,-4.0) .. (.95,-4.25) .. controls (.97,-4.6) and (1,-5.7) .. (1,-6.55) .. controls (1.04,-9.5) and (.99,-10.5) .. (1.015,-13) .. controls (1.045,-17.8) and (1.005,-23.5) .. (1.02,-30);}
\diagram{\begin{tikzcd}[ampersand replacement=\&,column sep=6em,row sep=4.5em]
\substack{\infty\\(x',x'')\deleted{\unsure{x}}\\P^1} \arrow[r,leftrightarrow,"u"{below},"{u^*(t)=\frac{t-x'}{t-x''}}"] \arrow[d,"f"] \& \substack{1\\(0,\infty)\\P^1} \arrow[d,"{h^*(t)=-t^3}"] \\
\substack{P^1\\(1,-1)\\\infty} \arrow[r,"v","{v^*(t)=\frac{1-t}{1+t}}"'] \& \substack{P^1\\(0,\infty)\\-1} \arrow[l,shift left=3.5,"{v^{-1*}(t)=\frac{1-t}{1+t})}"]
\end{tikzcd}}

\[
\begin{aligned}
f&=v^{-1}hu\\
f^*(t)&=u^*h^*v^{-1*}(t)=u^*h^*\left(\frac{1-t}{1+t}\right)\\
&=u^*\frac{1-h^*t}{1+h^*t}\mathbin{\unsure{=|:}}u^*\frac{1+t^3}{1-t^3}\\
&=\frac{1+u^*(t)^3}{1-u^*(t)^3}=
\end{aligned}
\]
\end{minipage}\hfill\begin{minipage}[t]{.47\linewidth}\scriptsize
Or ($x'$ et $x''$ étant racines pri\uline{mitives} $6^{\unsure{\text{èmes}|\text{ièmes}}}$ de 1)
\[
\begin{aligned}
{x'}^{3}={x''}^{3}&=-1,\quad\text{d'où}\quad {x''}^{3}-{x'}^{3}=0\\
\left\{\begin{array}{l}x'+x''=1\\x'x''=1\end{array}\right.&\\
{x'}^{2}+{x''}^{2}&=(x'+x'')^2-2x'x''=-1\\
(x'-x'')^2&=(x'+x'')^2-4x'x''=-3\\
x'-x''&=\sqrt{-3}\\
{x''}^{2}-{x'}^{2}&=(x''-x')(x''+x')=-\sqrt{-3}\\
\deleted{{x''}^{3}-{x'}^{3}=(x''-x')({x''}^{2}+x'x''+{x''}^{\unsure{2}})=\unsure{\sqrt{-3}(|-\sqrt{-3}(}}&
\end{aligned}
\]
\[
\begin{aligned}
f(t)&=\frac{1}{3\sqrt{-3}}\frac{2t^3-3t^2-3t+2}{t(t-1)}\\
&=\frac{1}{3\sqrt{-3}}\frac{(t+1)(2t-1)(t-2)}{t(t-1)}
\end{aligned}
\]
\end{minipage}
% archive sheet number: 37
